\subsection{从双流体逐项退到Hall、再到电阻和理想MHD}
现在把近似次序明确列出。第一步取准中性并使用低频Ampère关系；第二步在总动量中令$m_e/m_i\to0$，于是$\rho\simeq m_in$、$\vel\simeq\vel_i$，漂移应力按所保留尺度忽略；第三步在电子方程删掉惯性，但保留$\vel_i-\vel_e=\J/(en)$。这三步给
\begin{equation}
 \E=-\vel\times\B+{\J\times\B\over en}
      -{\grad p_e\over en}+\eta\J.
\end{equation}
这里“忽略电子惯性”是对动量项作排序，并非将每个含$m_e$的物理系数都机械置零；有限电阻的保留需要碰撞或有效非理想项的相应排序。代回Faraday定律，最外面的负号使四项依次变为感应方程中的
$+\vel\times\B$、$-\J\times\B/(en)$、$+\grad p_e/(en)$和$-\eta\J$。

若尺度进一步远大于$d_i$，且电子压力项在所需阶次也可忽略或其旋度为零，就得到电阻MHD：
\begin{equation}
 \partial_t\B=\grad\times(\vel\times\B)-\grad\times(\eta\J).
\end{equation}
只有$\eta$在空间均匀时，才可用$\grad\cdot\B=0$写成
\begin{equation}
 -\grad\times(\eta\J)
 =-{\eta\over\mu_0}\grad\times(\grad\times\B)
 =D_\eta\nabla^2\B.
\end{equation}
变电阻时还多出$-\grad\eta\times\J$。再要求所研究尺度上磁扩散可忽略，才有理想MHD。理想极限对薄电流层可能不一致：即使整体$S$很大，层内$\nabla^2\B$也可以很大。

\subsection{闭合一个Hall算例：还缺哪些方程}
以两温度、标量压力的Hall模型为例，已知$n,\vel,\B$并不等于已知$p_i,p_e$。一种无热流、无碰撞热交换的绝热闭合为
\begin{align}
 (\partial_t+\vel_i\cdot\grad)p_i+\gamma_i p_i\grad\cdot\vel_i&=0,\\
 (\partial_t+\vel_e\cdot\grad)p_e+\gamma_e p_e\grad\cdot\vel_e&=0,\qquad
 \vel_e=\vel_i-{\J\over en}.
\end{align}
两个压强被不同流速输运。在$\grad\cdot\J=0$时，离子和电子连续性与同一个$n$相容；但变密度下两种流的散度不必相同。若把电子压强也改由$\vel_i$输运，就额外作了近似，必须说明其尺度依据。

若改为$p_e=nT_{e0}$，便不再独立推进电子内能。等温是热力学闭合，不是从“电子质量小”自动推出来的。强压缩时，维持恒温意味着需要热输运或与热库交换；不能又要求这套等温方程满足孤立绝热能量守恒。

\subsection{无量纲化跟算：先变导数，再变系数}
写$\vct x=L_0\tilde{\vct x}$、$t=t_0\tilde t$、$\B=B_0\tilde\B$、$\vel=\VA\tilde\vel$、$n=n_0\tilde n$。于是
\begin{equation}
 \grad=L_0^{-1}\tilde\grad,\quad
 \partial_t={\VA\over L_0}\partial_{\tilde t},\quad
 \J={B_0\over\mu_0L_0}\tilde{\J},\quad
 \mathsf P={B_0^2\over\mu_0}\tilde{\mathsf P}.
\end{equation}
例如电子压力电场除以$E_0=\VA B_0$：
\begin{equation}
 {(\grad\cdot\mathsf P_e)/(en)\over\VA B_0}
 ={B_0\over\mu_0en_0L_0\VA}
 {\tilde\grad\cdot\tilde{\mathsf P}_e\over\tilde n}
 ={d_i\over L_0}
 {\tilde\grad\cdot\tilde{\mathsf P}_e\over\tilde n}.
\end{equation}
Hall和压力项在这一压强单位下带相同系数；若改用$n_0T_{e0}$作压力单位，压力系数还会带上$\beta_e/2$。两种写法可以相同，前提是压强变量的单位也随之改变。

电子惯性则有
\begin{equation}
 {(m_e/e)(\VA^2/L_0)\over\VA B_0}
 ={m_e\VA\over eB_0L_0}
 ={m_e\over m_i}{d_i\over L_0},
\end{equation}
其中用了$\VA=\Omega_{ci}d_i$。它与把电流代入后出现的$d_e^2/L_0^2$不矛盾，后者还使用了电子漂移速度及其空间梯度的标度。

\subsection{一个完整数值例：先用尺度判断该学哪套模型}
以下是专为练习选取的均匀上游参数，不对应某篇论文的测量：
\begin{equation}
 n_0=1\,{\rm cm^{-3}}=10^6\,{\rm m^{-3}},\quad
 B_0=20\,{\rm nT},\quad T_e=100\,{\rm eV},\quad
 L_0=1000\,{\rm km},
\end{equation}
离子为质子，质量比取真实值约1836。代入SI公式可得
\begin{align}
 \VA&\simeq436\,{\rm km\,s^{-1}},&
 t_0&\simeq2.29\,{\rm s},\\
 d_i&\simeq228\,{\rm km},&
 d_e&\simeq5.31\,{\rm km},\\
 \lambda_{De}&\simeq74.3\,{\rm m},&
 \rho_{\rm sound}&\simeq51.1\,{\rm km},\\
 \beta_e&\simeq0.101,&
 (\VA/c)^2&\simeq2.12\times10^{-6}.
\end{align}
所以在$L_0$尺度，电荷分离与位移电流的上述估计都很小；但$d_i/L_0\simeq0.228$，不能认为Hall作用必然微弱。若另取$S=10^4$，单层SP估计$\delta_{\rm SP}=10\,{\rm km}$，已远小于$d_i$，说明纯电阻MHD层的假设需要重新检查。

这里$B_0$就是所定义的上游场。若加上远大于它的导引场，回旋频率和声回旋半径应按相应磁化场重新计算，不能把同一张数值表直接沿用。该例训练的是“先查排序，再选择方程”，不是仅凭$d_i$一个数就完成模型选择。
